\chapter{Phân tích và thiết kế hệ thống}

Chương này cụ thể hóa nền tảng lý thuyết thành thiết kế hệ thống. Nội dung gồm yêu cầu chức năng và phi chức năng, kiến trúc tổng thể, thiết kế dữ liệu, lớp trừu tượng cho các phụ thuộc AI và hai luồng xử lý AI chính.

\medskip

\section{Phân tích yêu cầu}

Hệ thống cần đáp ứng đồng thời hai nhóm yêu cầu: nghiệp vụ thương mại điện tử và AI cá nhân hóa. Nhóm thứ nhất bao gồm duyệt sản phẩm, xem ảnh, quản lý giỏ hàng, thanh toán mô phỏng và thêm sản phẩm tại giao diện quản trị. Nhóm thứ hai bao gồm thiết lập hồ sơ cơ thể, quản lý tủ đồ cá nhân, tư vấn phối đồ bằng ngôn ngữ tự nhiên, thay thế sản phẩm tương tự và thử đồ ảo.

\medskip

Về mặt nghiệp vụ, các yêu cầu này liên kết thành một hành trình thống nhất. Người dùng có thể bắt đầu từ trang sản phẩm, thử đồ rồi thêm sản phẩm vào giỏ hàng; hoặc bắt đầu từ AI Stylist, nhận phương án phối đồ, thay một sản phẩm chưa phù hợp, thử toàn bộ trang phục và chuyển các sản phẩm được chọn sang giỏ hàng. Vì vậy, thiết kế phải liên kết được dữ liệu sản phẩm, phiên tư vấn, tủ đồ cá nhân và tác vụ VTON để giao diện trình bày trải nghiệm liền mạch.

\medskip

Hệ thống cũng phải phân biệt rõ dữ liệu nguồn và dữ liệu suy diễn. Sản phẩm, biến thể, giá và danh mục là dữ liệu nguồn do Medusa quản lý. Các thuộc tính thời trang như phong cách, hoàn cảnh sử dụng, màu sắc và chất liệu có thể được Gemini phân tích và ghi vào LanceDB để phục vụ tìm kiếm, nhưng không thay thế dữ liệu gốc. Tương tự, kết quả AI Stylist được lưu theo phiên để phục vụ lịch sử và truy vết, nhưng không làm thay đổi trực tiếp danh mục sản phẩm.

\medskip

Từ các yêu cầu nghiệp vụ trên, Bảng dưới đây tóm tắt những trường hợp sử dụng chính của hệ thống.

\subsection{Yêu cầu chức năng}

\begin{table}[H]
\centering
\small
\caption{Các trường hợp sử dụng chính của hệ thống}
\begin{tabularx}{\textwidth}{L{2.7cm}L{3.7cm}Y}
\toprule
\textbf{Nhóm chức năng} & \textbf{Trường hợp sử dụng} & \textbf{Mô tả} \\
\midrule
Cửa hàng & Xem và tìm kiếm sản phẩm & Người dùng xem danh sách sản phẩm từ Medusa Store API, tìm theo từ khóa và sắp xếp theo giá hoặc tên. \\
\addlinespace
Giỏ hàng & Thêm sản phẩm vào giỏ & Giao diện lưu giỏ hàng bằng Zustand, hỗ trợ thay đổi số lượng và mở bảng giỏ hàng. \\
\addlinespace
Hồ sơ AI & Thiết lập hồ sơ AI & Người dùng tải ảnh cơ thể, nhập chiều cao, cân nặng và giới tính; dữ liệu được lưu vào bảng \code{ai_profile}. \\
\addlinespace
Tủ đồ & Quản lý tủ đồ cá nhân & Người dùng tải ảnh trang phục, chọn loại trang phục, xóa trang phục và lưu kết quả VTON vào bộ sưu tập. \\
\addlinespace
AI Stylist & Truy vấn phối đồ & Người dùng nhập yêu cầu bằng tiếng Việt; hệ thống trả về phương án phối đồ, phần giải thích và danh sách sản phẩm. \\
\addlinespace
AI Stylist & Đổi sản phẩm tương tự & Người dùng chọn một sản phẩm trong phương án phối đồ để tìm sản phẩm tương đồng, đồng thời loại trừ các sản phẩm đang sử dụng. \\
\addlinespace
Virtual Try-On & Thử một sản phẩm & Người dùng thử trực tiếp một sản phẩm trong lưới sản phẩm trên ảnh cơ thể đã lưu. \\
\addlinespace
Virtual Try-On & Thử bộ trang phục & Người dùng thử bộ trang phục từ AI Stylist hoặc tủ đồ cá nhân. \\
\addlinespace
Thanh toán & Thanh toán mô phỏng & Thực hiện quy trình giữ tồn kho, tạo đơn hàng chờ xử lý, xử lý thanh toán và hoàn tất đơn hàng. \\
\addlinespace
Công cụ quản trị & Thêm sản phẩm & Giao diện \code{/admin-tools/add-product} tạo sản phẩm mới và kích hoạt đồng bộ vector. \\
\bottomrule
\end{tabularx}
\end{table}

\begin{figure}[H]
\centering
\resizebox{0.9\textwidth}{!}{%
\begin{tikzpicture}
    \node[diagram user, minimum width=2.5cm] (customer) at (0,0) {Khách hàng};
    \node[diagram user, minimum width=2.5cm] (admin) at (0,-6) {Người quản trị};

    \node[diagram small box, ellipse, minimum width=4.0cm] (browse) at (5,3) {Xem, tìm kiếm sản phẩm};
    \node[diagram small box, ellipse, minimum width=4.0cm] (profile) at (5,1) {Thiết lập AI Profile};
    \node[diagram small box, ellipse, minimum width=4.0cm] (wardrobe) at (5,-1) {Quản lý Wardrobe};
    \node[diagram small box, ellipse, minimum width=4.0cm] (stylist) at (5,-3) {Nhận tư vấn AI Stylist};
    \node[diagram small box, ellipse, minimum width=4.0cm] (adminproduct) at (5,-6) {Thêm sản phẩm};

    \node[diagram small box, ellipse, minimum width=4.0cm] (cart) at (10,3) {Quản lý giỏ hàng};
    \node[diagram small box, ellipse, minimum width=4.0cm] (vton) at (10,-0.5) {Thử đồ ảo VTON};
    \node[diagram small box, ellipse, minimum width=4.0cm] (replace) at (10,-3) {Đổi item tương tự};
    \node[diagram small box, ellipse, minimum width=4.0cm] (checkout) at (15,3) {Checkout MVP};

    \begin{scope}[on background layer]
        \node[diagram group, fill=blue!2, fit=(browse)(profile)(wardrobe)(stylist)(adminproduct)(cart)(vton)(replace)(checkout),
              label={[font=\small\bfseries]above:Hệ thống AI Fashion E-commerce}] {};
    \end{scope}

    \foreach \target in {browse,profile,wardrobe,stylist}
        \draw[line width=0.65pt] (customer) -- (\target.west);
    \draw[line width=0.55pt] (admin) -- (adminproduct.west);

    \draw[diagram logical] (browse) -- (cart);
    \draw[diagram logical] (cart) -- (checkout);
    \draw[diagram logical] (stylist) -- (replace);
    \draw[diagram logical] (profile) -- (vton);
    \draw[diagram logical] (wardrobe) -- (vton);
    \draw[diagram logical] (stylist) -- (vton);
    \draw[diagram logical] (wardrobe) -- (stylist);
\end{tikzpicture}%
}
\caption{Sơ đồ use case tổng quát của hệ thống}
\label{fig:use-case}
\end{figure}


\medskip

Các trường hợp sử dụng mô tả hành vi nhìn từ phía người dùng và hệ thống. Khả năng vận hành ổn định còn phụ thuộc vào các yêu cầu phi chức năng về độ trễ, xử lý bất đồng bộ, bảo mật nội bộ và tính nhất quán dữ liệu.

\subsection{Yêu cầu phi chức năng}

Hệ thống cần phân tách rõ các dịch vụ để tác vụ AI có tải tính toán lớn không làm gián đoạn luồng thương mại điện tử. Hàng đợi phải bảo đảm lưu giữ thông điệp, đồng thời mỗi bộ xử lý chỉ tiếp nhận một tác vụ tại một thời điểm để hạn chế quá tải GPU; số thông điệp chưa xác nhận được kiểm soát bằng cơ chế consumer prefetch của RabbitMQ \cite{rabbitmq-prefetch}. API nội bộ giữa Medusa và dịch vụ AI sử dụng mã truy cập nội bộ. Các phản hồi lỗi phải có cấu trúc nhất quán gồm trạng thái, mã lỗi và thông điệp.

\medskip

Giao diện cần phản hồi rõ các trạng thái chờ xử lý, đang xử lý, hoàn tất và lỗi trong quá trình VTON. Ảnh người dùng, bộ xử lý VTON và phương án phối đồ đang chọn được lưu bằng Zustand để tránh mất dữ liệu khi chuyển trang.

\medskip

Các thành phần AI áp dụng mô hình nhất quán cuối cùng. Khi sản phẩm được tạo hoặc cập nhật, PostgreSQL là nguồn dữ liệu chính; vector trong LanceDB được đồng bộ sau qua RabbitMQ \cite{lancedb,rabbitmq,postgresql}. Vì vậy, trong một khoảng thời gian ngắn, sản phẩm mới có thể đã xuất hiện trên cửa hàng nhưng chưa có trong kết quả AI Stylist. Đây là đánh đổi chấp nhận được vì luồng mua sắm vẫn hoạt động trong khi bộ tiêu thụ xử lý thông điệp đồng bộ.

\medskip

Về bảo mật, API nội bộ giữa Medusa và dịch vụ AI sử dụng tiêu đề \mbox{\texttt{x-internal-token}} để hạn chế truy cập trực tiếp từ máy khách. Thiết kế cho môi trường vận hành thực tế cần kết hợp xác thực phiên người dùng, phân quyền quản trị, giới hạn kích thước tệp tải lên và kiểm tra loại nội dung ảnh.

\medskip

Kiến trúc tổng thể dưới đây thể hiện cách giao diện người dùng, hệ thống nghiệp vụ thương mại điện tử, dịch vụ AI và các hệ lưu trữ phối hợp với nhau.

\section{Kiến trúc tổng thể}

Kho mã nguồn được tổ chức theo mô hình hợp nhất. Giao diện nằm trong \code{apps/web}, hệ thống nghiệp vụ nằm trong \code{apps/api-medusa}, dịch vụ AI nằm trong \code{apps/ai-service} và các lược đồ dữ liệu dùng chung được đặt trong \code{packages/shared-types}.

\medskip

Hình~\ref{fig:architecture-overview} thể hiện đường đi chính của dữ liệu giữa các thành phần.

\begin{figure}[H]
\centering
\resizebox{\textwidth}{!}{%
\begin{tikzpicture}
    \node[diagram user, minimum width=2.5cm, minimum height=1.3cm] (browser) at (0,0) {Trình duyệt\\người dùng};
    \node[diagram service, minimum width=3.4cm, minimum height=1.5cm] (next) at (4.0,0) {\textbf{Next.js Frontend}\\Storefront, AI Stylist,\\Wardrobe, VTON modal};
    \node[diagram service, minimum width=4.0cm, minimum height=1.7cm] (medusa) at (9.3,0) {\textbf{MedusaJS Backend / BFF}\\Custom API, workflow,\\AI Personalization, SSE};
    \node[diagram ai, minimum width=4.0cm, minimum height=1.7cm] (fastapi) at (15.3,0) {\textbf{FastAPI AI Service}\\Gemini adapter, CLIP,\\vector search, VTON worker};

    \node[diagram ai, minimum width=3.0cm] (gemini) at (13.2,3.5) {Google Gemini};
    \node[diagram ai, minimum width=3.4cm] (engines) at (18.2,3.5) {CatVTON local\\FLUX.2 / Replicate};

    \node[diagram data, minimum width=3.1cm] (postgres) at (5.8,-3.6) {PostgreSQL\\commerce + AI data};
    \node[diagram data, minimum width=2.5cm] (redis) at (9.5,-3.6) {Redis};
    \node[diagram queue, minimum width=3.2cm] (rabbit) at (13.2,-3.6) {RabbitMQ\\4 durable queues};
    \node[diagram data, minimum width=3.2cm] (lance) at (17.2,-3.6) {LanceDB\\product + closet vectors};
    \node[diagram data, minimum width=3.1cm] (static) at (21.3,-3.6) {Local static files\\VTON outputs};

    \draw[diagram arrow, <->] (browser) -- (next);
    \draw[diagram arrow, <->] (next) -- (medusa);
    \draw[diagram async] (medusa.south west) to[out=-150,in=-30] (next.south east);
    \draw[diagram arrow, <->] (medusa) -- (fastapi);

    \draw[diagram arrow, <->] (medusa) -- (postgres);
    \draw[diagram arrow, <->] (medusa) -- (redis);
    \draw[diagram async, <->] (medusa) -- (rabbit);
    \draw[diagram async, <->] (rabbit) -- (fastapi);
    \draw[diagram arrow, <->] (fastapi) -- (lance);
    \draw[diagram arrow, <->] (fastapi) -- (static);
    \draw[diagram arrow, <->] (fastapi) -- (gemini);
    \draw[diagram arrow, <->] (fastapi) -- (engines);

    \node[diagram note, below=9mm of rabbit, minimum width=5.4cm] (queues) {
        \code{product_metadata_sync}, \code{closet_metadata_sync}\\
        \code{ai_vision_jobs}, \code{ai_vision_results}
    };
    \draw[diagram logical] (rabbit) -- (queues);
\end{tikzpicture}%
}
\caption{Kiến trúc tổng thể của hệ thống theo các thành phần triển khai}
\label{fig:architecture-overview}
\end{figure}


Trong kiến trúc này, MedusaJS đóng vai trò hệ thống nghiệp vụ cốt lõi và BFF cho giao diện người dùng. Giao diện không gọi trực tiếp các tác vụ AI có tải tính toán lớn mà thông qua Medusa để bảo đảm dữ liệu sản phẩm, phiên tư vấn và tác vụ được lưu trong cơ sở dữ liệu. Dịch vụ AI tập trung vào xử lý ngôn ngữ, tìm kiếm vector và suy luận ảnh.

\medskip

Vai trò BFF của Medusa đặc biệt quan trọng vì kết quả từ dịch vụ AI chưa đủ để giao diện sử dụng trực tiếp. Dịch vụ AI có thể trả về định danh sản phẩm, định danh trang phục cá nhân, phần giải thích và danh sách mục được chọn; Medusa tiếp tục bổ sung tiêu đề, ảnh đại diện, biến thể và trạng thái tồn kho. Cách tổ chức này tập trung việc chuẩn hóa phản hồi và kiểm soát nghiệp vụ tại một thành phần.

\medskip

Kiến trúc tách rõ đường xử lý đồng bộ và bất đồng bộ. Các yêu cầu như xem sản phẩm, lưu hồ sơ AI hoặc gửi truy vấn AI Stylist sử dụng cơ chế yêu cầu--phản hồi HTTP. Ngược lại, VTON được đưa qua hàng đợi vì sinh ảnh có độ trễ cao và phụ thuộc tài nguyên tính toán.

\medskip

Từ kiến trúc tổng thể, thiết kế dữ liệu xác định cách tổ chức trạng thái nghiệp vụ, dữ liệu người dùng và dữ liệu truy hồi trước khi mô tả chi tiết các luồng AI.

\section{Thiết kế dữ liệu}

\subsection{Cơ sở dữ liệu quan hệ PostgreSQL}

MedusaJS quản lý dữ liệu thương mại điện tử cốt lõi như sản phẩm, biến thể, giá, danh mục và kênh bán hàng. Dữ liệu phục vụ cá nhân hóa được tách trong mô-đun \mbox{\texttt{ai-personalization}}, gồm năm thực thể chính.

\medskip

Bảng~\ref{tab:ai-personalization-entities} mô tả trách nhiệm và quan hệ chính của từng thực thể. Danh sách trường chi tiết được thể hiện trong mô hình dữ liệu và mã nguồn chuyển đổi lược đồ.

\begin{table}[H]
\centering
\small
\caption{Các thực thể của mô-đun cá nhân hóa AI}
\label{tab:ai-personalization-entities}
\begin{tabularx}{\textwidth}{L{3.4cm}L{5.0cm}Y}
\toprule
\textbf{Thực thể} & \textbf{Mục đích} & \textbf{Quan hệ chính} \\
\midrule
\code{ai_profile} & Lưu thông tin cơ thể và ảnh người dùng phục vụ cá nhân hóa và VTON. & Tham chiếu logic tới khách hàng Medusa theo \code{customer_id}. \\
\addlinespace
\code{user_closet_item} & Lưu trang phục cá nhân và kết quả VTON được người dùng giữ lại. & Mỗi khách hàng có thể sở hữu nhiều trang phục; dữ liệu được đồng bộ sang LanceDB. \\
\addlinespace
\code{stylist_session} & Lưu yêu cầu, lời giải thích và trạng thái của một phiên tư vấn. & Thuộc về một khách hàng và chứa nhiều mục gợi ý. \\
\addlinespace
\code{stylist_session_item} & Liên kết phiên tư vấn với sản phẩm được đề xuất hoặc sản phẩm thay thế. & Có khóa ngoại tới \code{stylist_session} và tham chiếu logic tới sản phẩm Medusa. \\
\addlinespace
\code{vton_job} & Lưu trạng thái và kết quả của tác vụ thử đồ ảo. & Các tác vụ nhiều bước được liên kết qua \code{parent_job_id}. \\
\bottomrule
\end{tabularx}
\end{table}

\begin{figure}[H]
\centering
\resizebox{0.98\textwidth}{!}{%
\begin{tikzpicture}[font=\scriptsize]
    \node[rectangle split, rectangle split parts=2, diagram box, align=left, minimum width=4.0cm] (customer) at (8,0) {
        \textbf{Medusa customer (core)}
        \nodepart{second}\code{id}
    };
    \node[rectangle split, rectangle split parts=2, diagram box, align=left, minimum width=4.0cm] (profile) at (2,-3.4) {
        \textbf{ai\_profile}
        \nodepart{second}\code{id}, \code{customer_id}\\
        \code{height}, \code{weight}, \code{gender}\\
        \code{base_body_image_url}
    };
    \node[rectangle split, rectangle split parts=2, diagram box, align=left, minimum width=4.0cm] (closet) at (8,-3.4) {
        \textbf{user\_closet\_item}
        \nodepart{second}\code{id}, \code{customer_id}\\
        \code{image_url}, \code{category}\\
        \code{metadata}, \code{vector_status}
    };
    \node[rectangle split, rectangle split parts=2, diagram box, align=left, minimum width=4.0cm] (product) at (14,-3.4) {
        \textbf{Medusa product (core)}
        \nodepart{second}\code{id}
    };
    \node[rectangle split, rectangle split parts=2, diagram box, align=left, minimum width=4.0cm] (session) at (2,-7.5) {
        \textbf{stylist\_session}
        \nodepart{second}\code{id}, \code{customer_id}\\
        \code{prompt}, \code{reasoning}, \code{status}
    };
    \node[rectangle split, rectangle split parts=2, diagram box, align=left, minimum width=4.0cm] (sessionitem) at (8,-7.5) {
        \textbf{stylist\_session\_item}
        \nodepart{second}\code{id}, \code{session_id}\\
        \code{product_id}, \code{replaces_item_id}
    };
    \node[rectangle split, rectangle split parts=2, diagram box, align=left, minimum width=4.6cm] (vton) at (14,-7.5) {
        \textbf{vton\_job}
        \nodepart{second}\code{id}, \code{customer_id}, \code{parent_job_id}\\
        \code{step}, \code{status}, \code{result_image_url}\\
        \code{error_message}, \code{retry_count}
    };

    \draw[diagram logical] (customer) -- (profile);
    \draw[diagram logical] (customer) -- (closet);
    \draw[diagram logical] (customer) -- (vton);
    \draw[diagram logical] (customer) -- (session);
    \draw[diagram arrow] (session) -- (sessionitem);
    \draw[diagram logical] (product) -- (sessionitem);
    \draw[diagram logical] (vton.south) to[out=-90,in=-15,looseness=4] (vton.south east);

    \node[diagram note, minimum width=9.2cm] (legend) at (8,-11.2) {
        Đường liền: quan hệ được khai báo bằng khóa ngoại. \quad
        Đường nét đứt: tham chiếu logic lưu dưới dạng ID.
    };
\end{tikzpicture}%
}
\caption{ERD của mô-đun cá nhân hóa AI và các tham chiếu tới dữ liệu lõi Medusa}
\label{fig:ai-erd}
\end{figure}


Hình~\ref{fig:ai-erd} phân biệt quan hệ khóa ngoại được khai báo trong mô-đun với các tham chiếu logic tới dữ liệu cốt lõi của Medusa. Cách phân biệt này giúp tránh giả định rằng mọi định danh liên hệ đều được ràng buộc bằng khóa ngoại vật lý.

\medskip

Mô hình dữ liệu quan hệ lưu những thông tin cần tính toàn vẹn và truy vấn theo người dùng. \code{ai_profile} có quan hệ logic một-một với khách hàng; \code{user_closet_item} có quan hệ một-nhiều; \code{stylist_session} và \code{stylist_session_item} liên kết phiên tư vấn với sản phẩm được đề xuất. \code{vton_job} có thể tạo chuỗi tác vụ thông qua \code{parent_job_id}.

\medskip

Cơ sở dữ liệu quan hệ không lưu trực tiếp vector nhúng. Vector được tách sang LanceDB vì loại dữ liệu này có phương pháp lập chỉ mục và truy vấn khác với dữ liệu giao dịch. PostgreSQL vẫn là nguồn chính cho trạng thái người dùng, phiên tư vấn và tác vụ; LanceDB là kho truy hồi phục vụ AI.

\subsection{Cơ sở dữ liệu vector LanceDB}

LanceDB đóng vai trò kho đọc chuyên biệt cho AI. Dữ liệu gốc vẫn nằm ở PostgreSQL và Medusa. Khi sản phẩm hoặc trang phục cá nhân thay đổi, Medusa phát sự kiện qua RabbitMQ để dịch vụ AI cập nhật vector tương ứng \cite{lancedb,rabbitmq}.

\medskip

Khác với PostgreSQL, LanceDB lưu vector nhúng và siêu dữ liệu phục vụ truy hồi ngữ nghĩa. Bảng sau trình bày hai bảng vector chính mà dịch vụ AI sử dụng.

\begin{longtable}{L{3.1cm}L{5.1cm}L{5.4cm}}
\caption{Lược đồ LanceDB phục vụ tìm kiếm ngữ nghĩa}
\label{tab:lancedb-schema} \\
\toprule
\textbf{Bảng} & \textbf{Nhóm dữ liệu} & \textbf{Nguồn đồng bộ} \\
\midrule
\endfirsthead
\toprule
\textbf{Bảng} & \textbf{Nhóm dữ liệu} & \textbf{Nguồn đồng bộ} \\
\midrule
\endhead
\code{products_vector} & Định danh sản phẩm, vector 512 chiều, thuộc tính thời trang, giá, tồn kho và dữ liệu hiển thị. & Sự kiện tạo hoặc cập nhật sản phẩm được phát vào \code{product_metadata_sync}. \\
\addlinespace
\code{closet_vector} & Định danh chủ sở hữu, vector 512 chiều, thuộc tính thời trang và địa chỉ ảnh. & Thao tác thêm hoặc xóa trang phục phát thông điệp vào \code{closet_metadata_sync}. \\
\bottomrule
\end{longtable}

Trong LanceDB, vector chỉ là một phần của bản ghi. Siêu dữ liệu đi kèm như danh mục, giới tính, màu sắc, giá, phong cách và trạng thái tồn kho cho phép hệ thống lọc ứng viên trước hoặc sau khi tìm kiếm tương đồng \cite{lancedb}. Nếu chỉ dựa vào độ tương đồng cosine, kết quả có thể gần về mặt ngữ nghĩa nhưng không phù hợp yêu cầu cụ thể. Bộ lọc thuộc tính giúp giới hạn không gian tìm kiếm và cung cấp thêm điều kiện để Medusa kiểm tra kết quả.

\medskip

Đối với tủ đồ cá nhân, trường \code{user_id} là ranh giới dữ liệu quan trọng. Vector của trang phục cá nhân chỉ được truy hồi trong phạm vi người dùng tương ứng. Trong phạm vi bản triển khai tối thiểu, ranh giới này phụ thuộc vào định danh khách hàng được truyền qua Medusa; khi vận hành thực tế, định danh phải được gắn với cơ chế xác thực đầy đủ và kiểm tra tại cả Medusa lẫn dịch vụ AI.

\medskip

Sau khi dữ liệu được tách thành các kho lưu trữ phù hợp, các phụ thuộc AI cần được cô lập để có thể thay mô hình hoặc nhà cung cấp mà không ảnh hưởng toàn bộ quy trình xử lý.

\section{Thiết kế lớp trừu tượng và bộ chuyển đổi}

Hệ thống sử dụng mẫu thiết kế Adapter để tránh phụ thuộc cứng vào một nhà cung cấp AI cụ thể. Ba giao diện chính lần lượt quy định khả năng phân tích ý định, sinh vector nhúng và thử đồ ảo. Các bộ chuyển đổi hiện tại triển khai Gemini, CLIP, CatVTON cục bộ và FLUX.2 đám mây.

\medskip

Cách thiết kế này cho phép thay đổi bộ xử lý AI trong tương lai. AI Stylist có thể chuyển từ Gemini sang nhà cung cấp mô hình ngôn ngữ khác mà không phải thay đổi toàn bộ tuyến API; VTON cũng có thể thay CatVTON cục bộ bằng một bộ xử lý đám mây khác.

\medskip

Mẫu thiết kế Adapter cho phép đánh giá thành phần theo hợp đồng giao tiếp thay vì phụ thuộc vào một thư viện cụ thể. \code{IIntentAnalyzer} quy định khả năng phân tích truy vấn và sinh phương án phối đồ có cấu trúc; \code{IEmbeddingService} bảo đảm vector có số chiều thống nhất; \code{IVirtualTryOnService} chuẩn hóa đầu vào ảnh và kết quả xử lý.

\medskip

Adapter không thay thế được kiểm thử tích hợp. Các mô hình AI có thể khác nhau về định dạng đầu ra, độ ổn định JSON, giới hạn kích thước ảnh, thời gian xử lý và cách báo lỗi. Do đó, ngoài kiểm thử riêng từng bộ chuyển đổi, hệ thống cần được kiểm tra theo luồng đầu cuối từ giao diện tới Medusa, RabbitMQ, dịch vụ AI và ngược lại.

\medskip

Trên nền dữ liệu và hợp đồng giao tiếp đã xác định, hai luồng AI Stylist và Virtual Try-On được thiết kế để xử lý lần lượt bài toán tư vấn phối đồ và sinh ảnh thử đồ.

\section{Thiết kế luồng AI Stylist}

Luồng AI Stylist bắt đầu khi người dùng nhập câu hỏi trong \code{ChatBox}. Giao diện gửi yêu cầu tới \code{/v1/stylist/search}. Medusa xác định \code{customer_id} từ ngữ cảnh xác thực, tiêu đề \mbox{\texttt{x-customer-id}} hoặc giá trị dự phòng \mbox{\texttt{mock-customer-id}}; thông tin giới tính trong hồ sơ AI được chuyển tới dịch vụ AI khi có.

\medskip

FastAPI tiếp nhận yêu cầu tại tuyến API nội bộ \code{/api/v1/stylist/search}. Tuyến này sử dụng \code{GeminiAnalyzerAdapter} để phân tích truy vấn. Nếu truy vấn không liên quan tới thời trang, Gemini trả về cờ \code{is_fashion_related} \code{=} \code{false} và hệ thống từ chối xử lý. Với truy vấn hợp lệ, hệ thống sử dụng \code{search_keywords} để gọi \code{VectorSearch.global_search}.

\medskip

Biểu đồ tuần tự dưới đây mô tả cách yêu cầu đi qua Medusa, FastAPI, Gemini và LanceDB trước khi giao diện nhận lại dữ liệu đã được bổ sung thông tin sản phẩm.

\begin{figure}[H]
\centering
\resizebox{\textwidth}{!}{%
\begin{tikzpicture}[x=2.75cm, y=-0.62cm, font=\scriptsize]
    \foreach \x/\name/\title in {
        0/fe/Frontend,
        1/medusa/Medusa BFF,
        2/fastapi/FastAPI,
        3/gemini/Gemini,
        4/lance/LanceDB,
        5/data/Product + AI Module
    } {
        \node[diagram small box, minimum width=2.25cm] (\name) at (\x,0) {\title};
        \draw[dashed, gray!65] (\name.south) -- ++(0,15.8);
    }

    \draw[diagram arrow] (fe.south |- 0,2.0) -- node[above, fill=white, inner sep=1pt] {1. \code{POST} \code{/v1/stylist/search}} (medusa.south |- 0,2.0);
    \draw[diagram arrow] (medusa.south |- 0,2.3) -- node[above, fill=white, inner sep=1pt] {2. lấy AI Profile / gender} (data.south |- 0,2.3);
    \draw[diagram arrow] (data.south |- 0,3.1) -- node[above, fill=white, inner sep=1pt] {profile hoặc \code{unisex}} (medusa.south |- 0,3.1);
    \draw[diagram arrow] (medusa.south |- 0,4.2) -- node[above, fill=white, inner sep=1pt] {3. internal \code{/api/v1/stylist/search}} (fastapi.south |- 0,4.2);
    \draw[diagram arrow] (fastapi.south |- 0,5.3) -- node[above, fill=white, inner sep=1pt] {4. \code{parse_query(query_text)}} (gemini.south |- 0,5.3);
    \draw[diagram arrow] (gemini.south |- 0,6.1) -- node[above, fill=white, inner sep=1pt] {intent, filters, keywords} (fastapi.south |- 0,6.1);

    \node[diagram note, minimum width=2.5cm] (guard) at (2,7.2) {Ngoài miền: trả refusal\\và kết thúc sớm};

    \draw[diagram arrow] (fastapi.south |- 0,8.5) -- node[above, fill=white, inner sep=1pt] {5. global search + user filter} (lance.south |- 0,8.5);
    \draw[diagram arrow] (lance.south |- 0,9.3) -- node[above, fill=white, inner sep=1pt] {pool store + closet} (fastapi.south |- 0,9.3);
    \draw[diagram arrow] (fastapi.south |- 0,10.4) -- node[above, fill=white, inner sep=1pt] {6. generate outfit từ pool} (gemini.south |- 0,10.4);
    \draw[diagram arrow] (gemini.south |- 0,11.2) -- node[above, fill=white, inner sep=1pt] {option IDs + reasoning} (fastapi.south |- 0,11.2);
    \draw[diagram arrow] (fastapi.south |- 0,12.3) -- node[above, fill=white, inner sep=1pt] {7. AI response} (medusa.south |- 0,12.3);
    \draw[diagram arrow] (medusa.south |- 0,13.4) -- node[above, fill=white, inner sep=1pt] {8. hydrate product / closet; lưu session} (data.south |- 0,13.4);
    \draw[diagram arrow] (data.south |- 0,14.2) -- node[above, fill=white, inner sep=1pt] {title, ảnh, variant, category} (medusa.south |- 0,14.2);
    \draw[diagram arrow] (medusa.south |- 0,15.3) -- node[above, fill=white, inner sep=1pt] {9. hydrated OutfitCard JSON} (fe.south |- 0,15.3);
\end{tikzpicture}%
}
\caption{Sequence diagram luồng AI Stylist từ frontend tới dữ liệu hiển thị}
\label{fig:stylist-sequence}
\end{figure}


\code{global_search} gộp hai nguồn dữ liệu: \code{products_vector} cho sản phẩm cửa hàng và \code{closet_vector} cho tủ đồ cá nhân. Kết quả được đánh dấu nguồn là \code{store} hoặc \code{closet}. Gemini sau đó sinh các phương án phối đồ, ưu tiên kết hợp trang phục cá nhân với sản phẩm cửa hàng khi phù hợp.

\medskip

Sau khi nhận dữ liệu từ dịch vụ AI, Medusa bổ sung dữ liệu giao dịch cho từng mục. Sản phẩm có định danh bắt đầu bằng \code{prod_} được truy vấn qua Product Module của Medusa; trang phục cá nhân được truy vấn qua AI Personalization Module. Medusa sau đó lưu \code{stylist_session}, các \code{stylist_session_item} hợp lệ và trả dữ liệu JSON cho giao diện.

\medskip

Việc lưu \code{stylist_session} phục vụ hai mục đích. Thứ nhất, phiên tư vấn cung cấp dữ liệu để hiển thị lịch sử và truy vết nội dung gợi ý. Thứ hai, định danh phiên tạo điểm nối để đo lường các hành động phát sinh từ gợi ý, chẳng hạn thêm sản phẩm vào giỏ hàng hoặc thực hiện thử đồ ảo.

\medskip

Khi người dùng yêu cầu đổi một sản phẩm, hệ thống không chạy lại toàn bộ quy trình gợi ý. Medusa truyền sản phẩm hiện tại và danh sách cần loại trừ tới dịch vụ AI để tìm kiếm vector trong ngữ cảnh tương đồng. Phần thay thế được xem là thao tác tinh chỉnh một phương án phối đồ, nhờ đó các sản phẩm còn lại và phần giải thích có thể được giữ nguyên.

\medskip

Khác với AI Stylist, luồng VTON không phù hợp với xử lý đồng bộ vì thời gian sinh ảnh dài và phụ thuộc tài nguyên tính toán. Thiết kế vì vậy tập trung vào hàng đợi, trạng thái tác vụ và kênh cập nhật kết quả về giao diện.

\section{Thiết kế luồng Virtual Try-On}

Luồng VTON được thiết kế bất đồng bộ. Giao diện mở kết nối SSE trước khi gửi yêu cầu tạo tác vụ \cite{sse-mdn}, sau đó gọi \code{/v1/vton/jobs} với ảnh người dùng, danh sách trang phục và bộ xử lý được chọn. Medusa tạo bản ghi \code{vton_job}, phát thông điệp vào hàng đợi \code{ai_vision_jobs} rồi trả về \code{job_id}.

\medskip

Dịch vụ AI tiêu thụ thông điệp từ \code{ai_vision_jobs}. Với chế độ cục bộ, bộ xử lý sử dụng \code{LocalCatVTONAdapter}; với chế độ đám mây, bộ xử lý sử dụng \code{CloudFluxVTONAdapter}. Kết quả được phát vào \code{ai_vision_results} sau khi xử lý hoàn tất.

\medskip

Bộ tiêu thụ của Medusa nhận \code{ai_vision_results}, gọi \code{VisionOrchestrator} để cập nhật cơ sở dữ liệu và kiểm tra tác vụ tiếp theo. Với VTON cục bộ nhiều bước, kết quả của bước trước trở thành ảnh đầu vào cho bước sau. Khi chuỗi xử lý hoàn tất hoặc thất bại, Medusa gửi dữ liệu tới giao diện qua \code{VisionSSEManager}.

\medskip

Luồng VTON nhiều bước được sử dụng khi người dùng thử một bộ gồm nhiều trang phục. CatVTON cục bộ thay lần lượt từng vùng trang phục, còn FLUX.2 đám mây có thể nhận nhiều ảnh tham chiếu trong một yêu cầu. Sự khác biệt được ẩn sau tham số \code{engine} để giao diện không phụ thuộc vào chi tiết cài đặt của từng bộ xử lý.

\medskip

Trạng thái của \code{vton_job} hỗ trợ chờ xử lý, đang xử lý, hoàn tất, thất bại và chuyển bước. Kết quả thành công được lưu tại \code{result_image_url}; lỗi được phản ánh bằng trạng thái \code{failed} và gửi tới giao diện qua SSE. Với chuỗi nhiều bước, \code{parent_job_id} và \code{step} xác định quan hệ giữa các tác vụ.

\medskip

Hình~\ref{fig:vton-flow} mô tả vòng đời của một tác vụ VTON, từ thời điểm giao diện tạo tác vụ đến khi kết quả được gửi về qua SSE.

\begin{figure}[H]
\centering
\resizebox{0.98\textwidth}{!}{%
\begin{tikzpicture}
    \node[diagram user, minimum width=3.5cm] (frontend) at (0,0) {Frontend mở SSE trước\\và gửi yêu cầu VTON};
    \node[diagram service, minimum width=3.7cm] (route) at (5.0,0) {Medusa \code{/v1/vton/jobs}\\kiểm tra payload và engine};
    \node[diagram service, minimum width=3.7cm] (jobs) at (10.2,0) {Tạo bản ghi \code{vton_job}\\trong AI Personalization};

    \node[diagram note, minimum width=3.8cm] (local) at (7.6,-2.5) {\textbf{Local}\\Tạo chuỗi job theo item;\\chỉ phát job đầu tiên};
    \node[diagram note, minimum width=3.8cm] (cloud) at (12.8,-2.5) {\textbf{Cloud}\\Tạo một job \code{outfit};\\gửi nhiều garment URL};

    \node[diagram queue, minimum width=3.4cm] (jobq) at (10.2,-5.2) {RabbitMQ\\\code{ai_vision_jobs}};
    \node[diagram ai, minimum width=3.8cm] (worker) at (15.2,-5.2) {FastAPI VTON consumer\\CatVTON hoặc FLUX.2};
    \node[diagram queue, minimum width=3.5cm] (resultq) at (20.2,-5.2) {RabbitMQ\\\code{ai_vision_results}};
    \node[diagram service, minimum width=3.8cm] (orchestrator) at (25.4,-5.2) {Medusa consumer +\\VisionOrchestrator};

    \node[diagram data, minimum width=3.2cm] (db) at (25.4,-1.8) {Cập nhật\\\code{vton_job}};
    \node[diagram data, minimum width=3.6cm] (output) at (15.2,-8.2) {Static output\\\code{/static/vton/*.png}};
    \node[diagram service, minimum width=3.6cm] (sse) at (25.4,-8.2) {VisionSSEManager\\broadcast theo user ID};
    \node[diagram note, minimum width=5.7cm] (childnote) at (9.0,-8.2) {
        Local còn child job: kết quả trước trở thành\\ảnh đầu vào mới cho bước tiếp theo
    };

    \draw[diagram arrow] (frontend) -- (route);
    \draw[diagram arrow] (route) -- (jobs);
    \draw[diagram arrow] (jobs) -- (local);
    \draw[diagram arrow] (jobs) -- (cloud);
    \draw[diagram async] (local) -- (jobq);
    \draw[diagram async] (cloud) -- (jobq);
    \draw[diagram async] (jobq) -- (worker);
    \draw[diagram arrow] (worker) -- (output);
    \draw[diagram async] (worker) -- (resultq);
    \draw[diagram async] (resultq) -- (orchestrator);
    \draw[diagram arrow] (orchestrator) -- (db);
    \draw[diagram async] (orchestrator) -- (sse);
    \draw[diagram async] (sse.south) |- (0,-10.5) -- (frontend.south);
\end{tikzpicture}%
}
\caption{Luồng Virtual Try-On bất đồng bộ qua RabbitMQ và SSE}
\label{fig:vton-flow}
\end{figure}


\medskip

Như vậy, thiết kế kết hợp xử lý đồng bộ cho AI Stylist với xử lý bất đồng bộ cho VTON, đồng thời duy trì ranh giới trách nhiệm giữa Medusa, dịch vụ AI, hàng đợi thông điệp và các hệ lưu trữ.

\clearpage
